\documentclass[10pt,a4paper]{amsart}
\usepackage[margin=19mm]{geometry}
\usepackage{amsmath,amssymb,mathtools}
\usepackage[scr=rsfs,cal=euler]{mathalfa}
\usepackage{xfrac}
\usepackage[unicode,hidelinks,bookmarks=false]{hyperref}
\newcommand{\ie}{i.e.\ }
\newcommand{\cf}{cf.\ }
\newcommand{\resp}{respectively\ }
\newcommand{\Subsection}{\subsection}
\providecommand{\nobreakdash}{\nobreak\hbox{-}}
\setlength{\emergencystretch}{2em}
\multlinegap0pt
\usepackage{mathtools}
\usepackage{amssymb}
\newtheorem{theorem}{Theorem}
\usepackage{cleveref}
\crefname{theorem}{Theorem}{Theorems}
\newcommand\cF{{\mathcal F}}
\title{The cohomological Kudla conjecture for unitary Shimura varieties}
\author{Fran\c{c}ois Greer, Salim Tayou}
\date{Source: arXiv:2507.13299v2 (CC BY 4.0)}
\begin{document}
\maketitle

\noindent\emph{Source and licence.} The cohomological Kudla conjecture for unitary Shimura varieties. Version arXiv:2507.13299v2 (CC BY 4.0), distributed by arXiv under the Creative Commons Attribution 4.0 International licence, http://creativecommons.org/licenses/by/4.0/, as recorded in the arXiv licence metadata for this exact version and shown on the abstract page https://arxiv.org/abs/2507.13299v2. Full licence notice: \texttt{licenses/arxiv-2507.13299-CC-BY-4.0.txt}.\par\smallskip

\noindent\textit{Editorial scope.} Counted unit: the article's theorem proposition-sl2 (source0.tex lines 696--699). The article's complete original proof of it is delivered with it (source0.tex lines 701--741, one \texttt{proof} environment of 41 lines). No mathematics is written, repaired or re-derived: the statement and the proof are the article's own lines, in the article's own notation, and no retained line is substituted or re-flowed.  Retained uncounted as the source-local context the proof needs: lines 613--617.  Nothing was omitted from any retained span, so the package carries no omission record.  No retained line contains a citation, so the wrapper carries no bibliography and no bibliography entry.  No retained line contains a figure environment, an image-inclusion command or drawing code, so no image is delivered and the package declares no asset dependency.  Editorial formatting only, in addition: an AMS \texttt{amsart} preamble that restates the article's own declarations and macros used by the retained lines, namely the environments \texttt{theorem} on separate counters, together with the standard packages those lines need.  The retained environments print in this document as Theorem 1 in order of appearance, whereas the article prints the same items with the numbering of its own full document.  The complete native source of the article as distributed in the arXiv e-print for this exact version is bundled unchanged as \texttt{sources/181765/source0.tex}.
\begin{align}\label{crucial}
   m_{ij}&=-Q(\lambda_1,\ldots,\widehat{\lambda}
   _i,\ldots,\underbrace{\lambda_i}_{(j-1)^{th}\, \textrm{position}},\ldots,\lambda_{g},\overline{\lambda}_1,\ldots,\widehat{\overline{\lambda}}_i,\ldots,\overline{\lambda}_{g})\\
   &=(-1)^{i+j}Q([\lambda]^g_j,[\overline{\lambda}]^g_i),
\end{align}

\begin{theorem}\label{proposition-sl2}
    The triple $(\Lambda,\Delta,H)$ is a $\mathfrak{sl}_2$-triple. In other words, the following relations hold:
    \[[\Lambda,\Delta]=H,\quad  [H,\Lambda]= 2\Lambda,\quad [H,\Delta]=-2\Delta~.\]
\end{theorem}

\begin{proof}
Let $P\in \cF_{n,g}$ and let $P_g=\Delta_g P$. We have: 
\begin{align}\label{eq:second-side}
\Lambda\circ\Delta_g (P)=\sum_{1\leq k,\ell\leq g}h(\lambda_k,\lambda_\ell)m_{\ell,k}(P_g)~.
\end{align}
On the other hand, we have:
 \begin{align*}
       \Delta_{g+1}\circ\Lambda( P)&= \Delta_{g+1}\left[\sum_{1\leq k,\ell\leq g+1}h(\lambda_k,\lambda_\ell)m_{\ell,k}(P)\right]\\
       &= \sum_{1\leq k,\ell\leq g+1}\Delta_{g+1}(h(\lambda_k,\lambda_\ell))m_{\ell,k}(P)+\sum_{1\leq k,\ell\leq g+1}h(\lambda_k,\lambda_\ell)\Delta_{g+1}(m_{\ell,k}(P))\\&+\sum_{1\leq k,\ell\leq g+1}\left[{\frac{^t\partial}{\partial \lambda_{g+1}}}h(\lambda_k,\lambda_\ell)\frac{\partial}{\partial \overline{\lambda}_{g+1}}m_{\ell,k}(P)+{\frac{^t\partial}{\partial \overline{\lambda}_{g+1}}}h(\lambda_k,\lambda_\ell)\frac{\partial}{\partial \lambda_{g+1}}m_{\ell,k}(P)\right]\\
       &= nP+\sum_{1\leq k,\ell\leq g+1}h(\lambda_k,\lambda_\ell)\Delta_{g+1}(m_{\ell,k}(P))\\&+\sum_{1\leq k,\ell\leq g+1}\left[{\frac{^t\partial}{\partial \lambda_{g+1}}}h(\lambda_k,\lambda_\ell)\frac{\partial}{\partial \overline{\lambda}_{g+1}}m_{\ell,k}(P)+{\frac{^t\partial}{\partial \overline{\lambda}_{g+1}}}h(\lambda_k,\lambda_\ell)\frac{\partial}{\partial \lambda_{g+1}}m_{\ell,k}(P)\right]~,\\
       \end{align*}
 In the last line and in what follows, we use the following simplifications: 
 \begin{enumerate}
     \item $m_{g+1,g+1}(P)=P$, $\Delta_{g+1}(h(\lambda_k,\lambda_\ell))=n$, if  $k=\ell=g+1$ and $0$ otherwise. 
     \item $\Delta_{g+1}(m_{g+1,k})=\Delta_{g+1}(m_{\ell,g+1})=0$ for all $1\leq \ell,k\leq g+1$, as $m_{g+1,k}$ and $m_{\ell,g+1}$ depend linearly on $\lambda_{g+1}$ and $\overline \lambda_{g+1}$ respectively.  
    \item For $\ell,k<g+1$, we have $\Delta_{g+1}(m_{\ell,k})=m_{\ell,k}(P_g)~,$ by commutation of derivatives.
    \item For $\ell,k<g+1$, using \cref{crucial}, one can check that:
    \[{\frac{^t\partial \left( h(\lambda_k,\lambda_\ell)\right)}{\partial \lambda_{g+1}}} \cdot \frac{\partial m_{\ell,k}(P)}{\partial \overline{\lambda}_{g+1}}=\delta_{k,g+1}  \, {^t\overline{\lambda}_\ell}\cdot \frac{\partial (m_{\ell,g+1}(P))}{\partial \overline{\lambda}_{g+1}}=-P\]
    and 
    \[{\frac{^t\partial(h(\lambda_k,\lambda_\ell))}{\partial \overline{\lambda}_{g+1}}}\cdot \frac{\partial(m_{\ell,k}(P))}{\partial \lambda_{g+1}}=\delta_{g+1,\ell}\, \lambda_k\cdot \frac{\partial (m_{g+1,k}(P))}{\partial \lambda_{g+1}}=-P~,\]
    \item If either $\ell=g+1$ or $k=g+1$ in the expressions in (4), then we get $0$, as $m_{g+1,g+1}(P)$ does not depend neither on $\lambda_{g+1}$, nor on $\overline{\lambda_{g+1}}$.
 \end{enumerate}

  Taking the previous relations into account, we get: 
  
       \begin{align*}
        \Delta_{g+1}\circ\Lambda( P)&= nP+\sum_{1\leq k,\ell\leq g}h(\lambda_k,\lambda_\ell)m_{\ell,k}(P_g)\\
        %&+\sum_{1\leq k\leq g+1}{\frac{^t\partial}{\partial \lambda_{g+1}}}h(\lambda_{g+1},\lambda_\ell)\frac{\partial}{\partial \overline{\lambda}_{g+1}}m_{\ell,g+1}(P)+\sum_{1\leq \ell\leq g+1}{\frac{^t\partial}{\partial \overline{\lambda}_{g+1}}}h(\lambda_k,\lambda_{g+1})\frac{\partial}{\partial \lambda_{g+1}}m_{g+1,k}(P)\\
        %&=nP+\sum_{1\leq k,\ell\leq g}h(\lambda_k,\lambda_\ell)\Delta_{g+1}(m_{\ell,k}(P))\\
        &+\sum_{1\leq \ell\leq g}\overline{\lambda}_\ell\cdot \frac{\partial}{\partial \overline{\lambda}_{g+1}}m_{\ell,g+1}(P)+\sum_{1\leq k\leq g}\lambda_k\cdot \frac{\partial}{\partial \lambda_{g+1}}m_{g+1,k}(P) \\
        &=(n-2g)P+\sum_{1\leq k,\ell\leq g}h(\lambda_k,\lambda_\ell)m_{\ell,k}(P_g)
   \end{align*}
 
Therefore,
\begin{align}\label{eq:one-side}
     \Delta_{g+1}\circ\Lambda(P)=(n-2g)P+\sum_{1\leq \ell,k\leq g}h(\lambda_i,\lambda_j)m_{i,j}(P_g)
\end{align}
   Finally, by subtracting \cref{eq:one-side} from \cref{eq:second-side}, we get
   \[[\Lambda,\Delta]P=(2g-n)P =H(P)~,\]
which proves the first relation. The last two relations are straightforward and we leave them to the reader.
\end{proof}

\end{document}
